%! Mean, Variance, and Covariance — Unit 8, Lesson 8.4 — Group Activity
\documentclass[10pt]{article}
\usepackage{linalg-article}
\usepackage{linalg-boxes}
\newcommand{\TallMath}[1]{$\displaystyle #1\rule[-1.4em]{0pt}{3.2em}$}

\begin{document}

\pageheader{Unit 8, Lesson 8.4}{Group Activity}
\namepartnerperiod

\vspace{0.10in}

\begin{headlinebox}{blush}
\small\textbf{The Shape of Data.} A data set is a cloud of points. Three numbers describe its shape: the
\textbf{mean} (center), the \textbf{variance} (spread $=$ average squared distance from the mean), and the
\textbf{covariance} (how two features move together --- $+$ together, $-$ oppositely, $0$ unrelated). Pack the
variances and covariance into the symmetric \textbf{covariance matrix} $V$. Work with your partner, show every
step, and \emph{explain} what you find.
\end{headlinebox}

\vspace{0.10in}

\begin{tcolorbox}[colback=white, colframe=black!40, title=\textbf{Tier R --- Remediate},
    fonttitle=\bfseries, arc=1.5mm, left=3mm, right=3mm, top=2mm, bottom=2mm, breakable]
A group of four students scored $1,3,5,7$ on a quiz.
\begin{enumerate}[leftmargin=1.7em, itemsep=8pt]
  \item \textbf{Mean.} Find the mean $\mu=\dfrac{1+3+5+7}{4}=\blank{1.0cm}$.
  \item \textbf{Variance.} Subtract the mean, square, and average:
        \[ \sigma^2=\frac{(1-\mu)^2+(3-\mu)^2+(5-\mu)^2+(7-\mu)^2}{4}=\blank{1.4cm}. \]
  \item \textbf{Standard deviation.} $\sigma=\sqrt{\sigma^2}=\blank{1.4cm}$ (leave as a square root).
  \item In one sentence, what would it mean for \emph{another} class to have a \emph{larger} variance? \writeline
\end{enumerate}
\end{tcolorbox}

\begin{tcolorbox}[colback=white, colframe=black!40, title=\textbf{Tier A --- Approaching Proficiency},
    fonttitle=\bfseries, arc=1.5mm, left=3mm, right=3mm, top=2mm, bottom=2mm, breakable]
Four students have two scores each (hours slept, then test score), giving points
$(2,7),(4,9),(6,3),(8,5)$.
\begin{enumerate}[leftmargin=1.7em, itemsep=9pt]
  \item \textbf{Means.} $\mu_x=\blank{0.9cm}$, \quad $\mu_y=\blank{0.9cm}$. Then center: list the deviations
        \[ x-\mu_x:\ \blank{2.6cm}, \qquad y-\mu_y:\ \blank{2.6cm}. \]
  \item \textbf{Variances.} $\sigma_x^2=\blank{1.0cm}$, \quad $\sigma_y^2=\blank{1.0cm}$.
  \item \textbf{Covariance.} $\sigma_{xy}=\dfrac{1}{4}\displaystyle\sum (x_i-\mu_x)(y_i-\mu_y)=\blank{1.2cm}$.
  \item \textbf{Build $V$.} Fill in the covariance matrix, and say what its \emph{sign} tells you about sleep and scores:
        \[ V=\begin{bmatrix}\ \blank{0.8cm}\ &\ \blank{0.8cm}\ \\[2pt] \blank{0.8cm}\ &\ \blank{0.8cm}\ \end{bmatrix}. \]
        \writeline
\end{enumerate}
\end{tcolorbox}

\begin{tcolorbox}[colback=white, colframe=black!40, title=\textbf{Tier E --- Extension},
    fonttitle=\bfseries, arc=1.5mm, left=3mm, right=3mm, top=2mm, bottom=2mm, breakable]
\textbf{Principal components of the covariance matrix.} Use the notes' matrix
$V=\left[\begin{smallmatrix}5&3\\3&5\end{smallmatrix}\right]$.
\begin{enumerate}[leftmargin=1.7em, itemsep=9pt]
  \item \textbf{Eigenvalues.} Solve $\det(V-\lambda I)=(5-\lambda)^2-9=0$: \quad $\lambda=\blank{0.9cm}$ or $\blank{0.9cm}$.
  \item \textbf{Eigenvectors (principal components).} For each $\lambda$, find a direction:
        \[ \lambda=8:\ \blank{1.8cm}, \qquad \lambda=2:\ \blank{1.8cm}. \]
        Which one is PC1 (the direction of greatest spread), and what fraction of the total variance does it hold? \writeline
  \item \textbf{Justify.} In one or two sentences each: (a) why is every covariance matrix \emph{symmetric}?
        (b) Why does the \emph{trace} of $V$ equal the total variance \emph{and} the sum of the eigenvalues? \writeline
\end{enumerate}
\end{tcolorbox}

\end{document}
